% ===================================================================
%  اللوحة نمرة ١٢ — المحطة. --- السكك الحديدية. --- المراكب.
%
%  Traduction, faite en 2026, en arabe egyptien ecrit, sur l'ido de
%  texto/io. Regles de la colonne : voir 10-lawha-01.tex. Une seule
%  numerotation.
%
%  LE CHEMIN DE FER PORTE LE SIXIEME METIER EN -جي DE LA COLONNE, ET
%  C'EST LE PLUS SPECIALISE DE TOUS :
%
%    تحويلجي     l'aiguilleur — celui qui « fait tourner »
%
%  Apres صفرجي (t4), عربجي et كهربجي (t5), جزمجي (t7) et بوسطجي
%  (t11), il n'y a plus a demontrer que le procede est vivant : il
%  s'applique meme a un metier ne avec la vapeur.
%
%  LA GARE EGYPTIENNE A SES TITRES ET SES BUREAUX :
%
%    كمساري (6)         le controleur   de l'italien commissario
%    ناظر المحطة (13)   le chef de gare — le meme ناظر qu'au tableau 1
%    وكيل الناظر (52)   le sous-chef
%    شيّال (16)          le porteur      — celui du tableau 5
%    مكتب الأمانات (25) la consigne
%    شبّاك (8)           le guichet
%    بوفيه (34)         le buffet
%    جرسون (35)         le garcon
%
%  ET LA VOIE S'APPELLE الشريط (63, 64), « le ruban ». Ce n'est pas
%  une image : c'est le mot ordinaire de la voie ferree en Egypte, et
%  le passage a niveau qui la coupe est المزلقان. Le francais dit
%  « les rails », le cantonais 路軌, l'egyptien nomme la ligne
%  entiere par sa forme.
%
%  ET « وابور » REVIENT UNE QUATRIEME FOIS (54), pour la locomotive.
%  Il avait fait la locomobile au tableau 8 et le paquebot au tableau
%  10. Un seul mot pour la vapeur, quel que soit ce qu'elle traine.
%
%  DEUX BLOCS SONT LES MEMES PIEGES QUE PARTOUT : le bureau du
%  telegraphe (t12-02-2), qui vient APRES le controleur a qui on le
%  demande et qui avait fait trebucher fr-CA et yue ; et la valise
%  avant la sortie (t12-06-1).
% ===================================================================

%%K t12-apar-1 apar
\VUsaut{4.0mm}
\VUtitre{15.0pt}{60}{اللوحة نمرة ١٢}
\VUsaut{2.4mm}
\VUpk{11.4pt}{40}{المحطة. --- السكك الحديدية. --- المراكب.}
\VUsaut{3.4mm}
\textsc{ملحوظة.} --- \textit{جداول المواعيد المستعملة في كل بلد مختلفة، علشان كده
إحنا بناخد كمثال مواعيد} \VUgras{الجدول الفرنساوي}\textit{.}

%%K t12-01-1 p
1. --- سافرت حالًا في ألمانيا. وقبل ما أسافر، اشتريت \VUgras{دليل القطارات}
\textsuperscript{(15)} علشان أعرف مواعيد قيام القطر. وبعد ما اتأكّدت إن أنسب قطر
هيقوم من باريس الساعة تسعة بالليل ويوصل ستراسبورج الساعة واحدة وتلاتاشر دقيقة،
حضّرت سفريتي، وتاني يوم خلّيتهم يوصّلوني المحطة.

%%K t12-02-1 p
2. --- ولما دخلت صالة القيام الكبيرة، ورا \VUgras{قسّيس} \textsuperscript{(2)} ريفي
عجوز شايل بإيده \VUgras{شنطة السفر} \textsuperscript{(3)} بتاعته، كان في
\VUgras{مسافرين} \textsuperscript{(1)} كتير وصلوا خلاص.

%%K t12-02-2 p
وعلشان كنت عايز أبعت \VUgras{تلغراف} \textsuperscript{(5)}، خلّيت \VUgras{الكمساري}
\textsuperscript{(6)} يوضّح لي \VUgras{مكتب التلغراف} \textsuperscript{(4)}.

%%K t12-03-1 p
3. --- وقبل ما أدخل \VUgras{صالة الانتظار} \textsuperscript{(7)} رحت أخد
\VUgras{تذكرة} \textsuperscript{(9)} رايح جاي.

%%K t12-04-1 p
4. --- وما استنّيتش كتير على \VUgras{الشبّاك} \textsuperscript{(8)}، علشان ما كانش
في ناس كتير هناك. و\VUgras{عسكري} \textsuperscript{(10)} كان مسافر في أجازة تكرّم
وسابلي دوره، وكده كان عندي وقت كفاية علشان أسأل \VUgras{ناظر المحطة}
\textsuperscript{(13)} كام سؤال، وهو كان بيتكلّم مع \VUgras{مأمور} \textsuperscript{(12)}
المراقبة ومع \VUgras{الجندرمي} \textsuperscript{(11)} اللي مناوب.

%%K t12-tit-1 sub
\VUpk{10.2pt}{40}{تسجيل العفش.}
\VUsaut{3.0mm}

%%K t12-05-1 p
5. --- وبعد ما اشتريت جرنال من \VUgras{كشك الجرايد} \textsuperscript{(14)}، خلّيتهم
يوزنوا \VUgras{شنطي} \textsuperscript{(17)} اللي \VUgras{الشيّال} \textsuperscript{(16)}
كان جايبها حالًا على \VUgras{عربية العفش} \textsuperscript{(19)}؛ و\VUgras{الموظّف}
\textsuperscript{(22)} المسؤول عن \VUgras{الميزان} \textsuperscript{(21)} قال لي إن شنطتي
وزنها خمسة وتلاتين كيلو؛ يعني \VUgras{زيادة} خمس كيلو، لازم أدفع عنها زيادة، على
\VUgras{شبّاك} العفش \textsuperscript{(23)}.

%%K t12-06-1 p
6. --- وعلشان كنت بدري أوي، كان عندي وقت فاضي علشان أتفرّج على حركة المسافرين
في المحطة. في ناس بتودّي أو تاخد شنطها الخفيفة من \VUgras{مكتب الأمانات}
\textsuperscript{(25)}؛ وناس تانية بتبصّ على \VUgras{جدول المواعيد}
\textsuperscript{(26)}. والمسافرين الواصلين بيستعجلوا على \VUgras{الخروج}
\textsuperscript{(32)} و\VUgras{الشنطة} \textsuperscript{(24)} في إيدهم. وكام واحد
بيستنّوا في \VUgras{مكتب تسليم العفش} \textsuperscript{(27)} وبيقدّموا
\VUgras{إيصالهم} \textsuperscript{(29)} علشان ياخدوا شنطهم، أو \VUgras{صناديقهم}
\textsuperscript{(28)}، أو \VUgras{مقاطفهم} \textsuperscript{(30)}.

%%K t12-07-1 p
7. --- و\VUgras{موظّفين المكوس} \textsuperscript{(33)} بيفتّشوا الطرود بعناية علشان
يشوفوا في حاجة لازم تتقال ولا لأ.

%%K t12-08-1 p
8. --- وفي مسافرين راحوا \VUgras{البوفيه} \textsuperscript{(34)}، اللي
\VUgras{الجرسون} \textsuperscript{(35)} بيخدمهم فيه بنشاط. ولازم يستعجلوا، علشان
\VUgras{القطر} \textsuperscript{(36)}، اللي هو إكسبريس، مش هيقعد كتير في المحطة.

%%K t12-09-1 p
9. --- وبعدين رحت رصيف القيام ودوّرت على \VUgras{عربية} \textsuperscript{(37)} مباشرة
علشان ما أضطرّش أغيّر العربية في السفر. وطلعت عربية فيها ممرّ وفيها
\VUgras{كابينة} \textsuperscript{(38)} للمدخّنين وكابينة محجوزة «للسيدات لوحدهن».

%%K t12-tit-2 sub
\VUpk{10.2pt}{40}{القيام بالليل.}
\VUsaut{3.0mm}

%%K t12-10-1 p
10. --- وبعد ما طلعت \VUgras{السلّمة} \textsuperscript{(40)}، وقفلت \VUgras{الباب}
\textsuperscript{(39)}، ولفّيت \VUgras{الستارة} \textsuperscript{(41)}، وحطّيت شنطي
الصغيرة على \VUgras{الرفّ} \textsuperscript{(43)}، قعدت على \VUgras{الكنبة}
\textsuperscript{(42)}. ولما شفت \VUgras{بتاع اللمبات} \textsuperscript{(44)} بيولّع
اللمبات، افتكرت إني هقضّي ليلة في العربية، وناديت الموظّف المسؤول عن تأجير
\VUgras{المخدّات} \textsuperscript{(45)} علشان أطلب واحدة.

%%K t12-11-1 p
11. --- وكمساري القطر جه يراجع التذاكر. وسألته ليه لسه ما قمناش، والمعاد جه
خلاص. ردّ عليّ إن \VUgras{رئيس القطر} \textsuperscript{(46)} مشغول بتحميل عجل في
\VUgras{عربية العفش} \textsuperscript{(47)}. وكمان لسه ما غيّروش كل
\VUgras{مدفّيات الرجلين} \textsuperscript{(48)}.

%%K t12-12-1 p
12. --- بس كنّا قربنا نقوم، علشان \VUgras{الوقّاد} \textsuperscript{(50)}، على
\VUgras{عربية الفحم} \textsuperscript{(49)} بتاعته، خد فحم علشان يشعّل نار
\VUgras{الوابور} \textsuperscript{(54)}، و\VUgras{السايق} \textsuperscript{(51)} كان
مستنّي بس إشارة \VUgras{وكيل الناظر} \textsuperscript{(52)}، اللي كان على
\VUgras{الرصيف} \textsuperscript{(53)}.

%%K t12-13-1 p
13. --- و\VUgras{غلاية} \textsuperscript{(56)} الوابور كانت بتهدر بصوت مكتوم؛
و\VUgras{المدخنة} \textsuperscript{(55)} كانت بتقذف سيول دخان مخلوطة
بـ\VUgras{البخار} \textsuperscript{(60)}. و\VUgras{صفّارة} \textsuperscript{(59)} حادّة
صوّتت بقوة و\VUgras{المكابس} \textsuperscript{(58)} ابتدت تتحرّك، وهي بتدفع
\VUgras{عجل} \textsuperscript{(57)} الماكينة التقيلة.

%%K t12-tit-3 sub
\VUsaut{3.0mm}
\VUpk{10.2pt}{40}{قيام!}
\VUsaut{3.0mm}

%%K t12-14-1 p
14. --- وسرعة القطر، وهو ماشي على \VUgras{الشريط} \textsuperscript{(64)}، زادت؛
و\VUgras{الأرصفة} \textsuperscript{(65)} اختفت؛ وعدّينا \VUgras{بيوت المية}
\textsuperscript{(66)}، وكمان عدّينا

%%K t12-noto-1 noto
\VUnotes{143.2mm}{%
\textit{(*) ملحوظة.} --- \textit{طبعًا وصف السفرية على حسب حدود ما قبل الحرب.}%
}

%%K t12-14-1 p suite
قطر واقف على \VUgras{الشريط} \textsuperscript{(63)} التاني: \VUgras{عربيات نوم}
\textsuperscript{(67)}، و\VUgras{عربية مطعم} \textsuperscript{(68)}، و\VUgras{عربية
بوستة} \textsuperscript{(69)}.

%%K t12-15-1 p
15. --- وبعدين، \VUgras{محطة البضايع} \textsuperscript{(71)} عدّت قدام عينينا. وتحت
الهناجر، موظّفين بيشحنوا \VUgras{البضايع} \textsuperscript{(72)} المرسلة بالقطر
البطيء. و\VUgras{عمال المناورة} \textsuperscript{(76)} جنب \VUgras{خزان المية}
\textsuperscript{(73)} بيناوروا \VUgras{عربيات المواشي} \textsuperscript{(75)}
و\VUgras{العربيات المسطّحة} \textsuperscript{(79)} علشان يعملوا \VUgras{قطر بضايع}
\textsuperscript{(80)}.

%%K t12-16-1 p
16. --- وعدّينا آخر \VUgras{مفرق} \textsuperscript{(78)}، اللي كان واقف جنبه
التحويلجي؛ وعدّينا \VUgras{قرص الإشارة} \textsuperscript{(86)} والمحطة اختفت بعيد.

%%K t12-17-1 p
17. --- وقريّب عدّينا \VUgras{كوبري حديد} \textsuperscript{(74)} على
\VUgras{النهر} \textsuperscript{(81)}. وبعد ما عدّينا الكوبري، مشينا جنب النهر،
اللي عليه \VUgras{لنشات} \textsuperscript{(82)} قوية بتجرّ \VUgras{صنادل}
\textsuperscript{(83)} تقيلة. وشفنا كمان \VUgras{باخرة ركّاب} \textsuperscript{(84)}
وهي بترسي على \VUgras{رصيف عايم} \textsuperscript{(85)}.

%%K t12-18-1 p
18. --- وبعد ما عدّينا بسرعة كبيرة كام محطة، عدّينا كام \VUgras{نفق}
\textsuperscript{(87)} وسبنا ورانا \VUgras{بيوت} \textsuperscript{(88)} صغيرة ما
تتعدّش لـ\VUgras{خفرا المزلقان}. وهدهدني هزّ القطر، فنمت نوم عميق.

%%K t12-tit-4 sub
\VUsaut{3.0mm}
\VUpk{10.2pt}{40}{محطة دويتش أفريكور.}
\VUsaut{3.0mm}

%%K t12-19-1 p
19. --- وصحيت فجأة على صوت موظّف بيزعّق: «دويتش أفريكور!
\textsuperscript{(*)}. كل الناس تنزل!» وحصل تفتيش الجمرك وكل الإجراءات المملّة
بتاعة الحدود. واضطرّيت أظبط ساعتي على \VUgras{الساعة الكبيرة} \textsuperscript{(89)}
بتاعة المحطة، اللي كانت سابقة بخمسة وخمسين دقيقة؛

%%K t12-19-1 p suite
وأنا لسه صاحي بالكاد، اتلخبطت بين \VUgras{العقرب الكبير} \textsuperscript{(90)}
و\VUgras{الصغير} \textsuperscript{(91)}؛ و\VUgras{وش الساعة} \textsuperscript{(92)} نفسه
كان مش باين خالص من شبورة الصبح.

%%K t12-20-1 p
20. --- ولما موظّفين الجمرك كانوا بيفتّشوا، فضلت أتثاءب وأقرا \VUgras{الإعلانات}
\textsuperscript{(93)} اللي مزيّنة حيطان المحطة. وفطرت علشان أقضّي الوقت، وبصّيت
على خريطة \VUgras{شبكة السكك} \textsuperscript{(94)} الألمانية علشان أشوف فاضل قدّ إيه
لستراسبورج، ورجعت عربيتي، وأنا متقوّي بالأمل إني قريّب هوصل لهدف سفريتي.

\VUsaut{6.0mm}
\VUfilet{22mm}
